To dissect the contribution of individual components to MPDenoiseNet's overall performance, we conducted an ablation study by systematically removing or reducing key architectural elements.  This analysis allows us to quantify the importance of each module and understand their specific roles in the denoising process.

\subsection{Ablated Model Configurations}

We evaluated several variants of the MPDenoiseNet architecture: \textbf{MPDenoiseNet (5 Transformer Blocks)}, which reduces the Transformer stack from 10 to 5; \textbf{MPDenoiseNet w/o Transformer Blocks}, removing the stack entirely; \textbf{MPDenoiseNet w/o Anisotropic Diffusion}, removing the AD branch; \textbf{MPDenoiseNet w/o SeConv}, removing the SeConv branch; and \textbf{MPDenoiseNet w/o AutoEncoder}, removing the final reconstruction module. All other architectural parameters and training procedures remained consistent with the full MPDenoiseNet model described in the previous sections, ensuring a controlled ablation study.

The ablation study followed the same experimental methodology as described in Section \ref{sec:experimental_results}. Below are the average PSNR and SSIM scores of each ablation experiment (average across all datasets). Based on the performance degradation observed when components were removed (Tables \ref{table:ablation-1}, \ref{table:ablation-2}, \ref{table:ablation-3}), the MPDenoiseNetomponents rank in order of impact as follows: Transformer Blocks, U-Net Module, SeConv Blocks, and finally Anisotropic Diffusion Blocks. 

While this ablation study indicates the Anisotropic Diffusion Blocks have the most modest quantitative impact among the model's components, it is included on the merit of being computationally inexpensive, while still providing a performance bonus. Further qualitative analysis on the denoised images is needed to exclude whether the AD branch provides a visual quality benefit, particularly in edge preservation.

\begin{table}[H]
    \centering
    \begin{tabular}{>{\centering\arraybackslash}p{2.5cm}*{10}{>{\centering\arraybackslash}p{.8cm}}}
        \toprule
        & \multicolumn{4}{c}{PSNR} & \multicolumn{4}{c}{SSIM} \\
        \cmidrule(lr){2-5} \cmidrule(lr){6-9}
        & 15 & 25 & 50 & 60 & 15 & 25 & 50 & 60 \\
        \midrule
        5 Transformers & 31.35 & 28.98 & 25.90 & 25.12 & 0.9040 & 0.8541 & 0.7543 & 0.7227 \\
        No Anisotropic & 31.26 & 28.93 & 25.85 & 25.07 & 0.9058 & 0.8539 & 0.7532 & 0.7221 \\
        No SeConv & 30.66 & 28.35 & 25.38 & 24.62 & 0.8992 & 0.8440 & 0.7389 & 0.7063 \\
        No Transformers & 29.22 & 25.96 & 24.80 & 24.47 & 0.8377 & 0.7337 & 0.7187 & 0.6985 \\
        No U-Net & 30.97 & 28.54 & 25.31 & 24.48 & 0.8924 & 0.8348 & 0.7222 & 0.6864 \\
        \underline{Baseline} & \underline{31.38} & \underline{29.08} & \underline{25.99} & \underline{25.19} & \underline{0.9057} & \underline{0.8566} & \underline{0.7560} & \underline{0.7242} \\
        \bottomrule
    \end{tabular}
    \caption{Ablation Study Results - Gaussian Test}
    \label{table:ablation-1}
\end{table}

\begin{table}[!hbt]
    \centering
    \begin{tabular}{>{\centering\arraybackslash}p{2.5cm}*{10}{>{\centering\arraybackslash}p{.55cm}}}
        \toprule
        Dataset & 150 & 160 & 170 & 180 & 190 & 200 & 210 & 220 & 230 & 240 \\
        \midrule
        \textbf{BERNOULLI} \\
        5 Transformers & 29.70 & 29.08 & 28.45 & 27.78 & 27.09 & 26.33 & 25.50 & 24.55 & 23.43 & 21.97 \\
        No Anisotropic & 29.59 & 29.00 & 28.36 & 27.71 & 27.02 & 26.27 & 25.43 & 24.49 & 23.39 & 21.94 \\
        No SeConv & 29.59 & 28.98 & 28.33 & 27.64 & 26.91 & 26.16 & 25.31 & 24.37 & 23.26 & 21.82 \\
        No Transformers & 28.52 & 27.95 & 27.36 & 26.74 & 26.08 & 25.37 & 24.60 & 23.71 & 22.68 & 21.32 \\
        No U-Net & 28.72 & 28.10 & 27.47 & 26.79 & 26.08 & 25.32 & 24.51 & 23.58 & 22.52 & 21.20 \\
        \underline{Baseline} & \underline{30.10} & \underline{29.48} & \underline{28.84} & \underline{28.15} & \underline{27.43} & \underline{26.64} & \underline{25.78} & \underline{24.81} & \underline{23.61} & \underline{22.06} \\
        \midrule
        \textbf{POISSON} \\
        5 Transformers & 30.79 & 30.59 & 30.44 & 30.17 & 30.04 & 29.79 & 29.63 & 29.53 & 29.40 & 29.22 \\
        No Anisotropic & 30.82 & 30.60 & 30.46 & 30.21 & 29.98 & 29.86 & 29.63 & 29.53 & 29.39 & 29.17 \\
        No SeConv & 30.65 & 30.40 & 30.26 & 30.08 & 29.87 & 29.73 & 29.54 & 29.38 & 29.16 & 29.00 \\
        No Transformers & 29.46 & 29.26 & 29.08 & 28.89 & 28.76 & 28.54 & 28.36 & 28.21 & 28.01 & 27.81 \\
        No U-Net & 30.38 & 30.09 & 29.91 & 29.69 & 29.50 & 29.31 & 29.10 & 28.94 & 28.72 & 28.56 \\
        \underline{Baseline} & \underline{31.48} & \underline{31.19} & \underline{30.99} & \underline{30.83} & \underline{30.58} & \underline{30.34} & \underline{30.16} & \underline{30.00} & \underline{29.79} & \underline{29.54} \\
        \midrule
        \textbf{SAP} \\
        5 Transformers & 29.24 & 28.65 & 28.03 & 27.42 & 26.75 & 26.00 & 25.17 & 24.26 & 23.16 & 21.71 \\
        No Anisotropic & 29.20 & 28.59 & 28.00 & 27.35 & 26.65 & 25.94 & 25.12 & 24.18 & 23.07 & 21.65 \\
        No SeConv & 28.67 & 28.20 & 27.68 & 27.05 & 26.19 & 25.42 & 24.68 & 23.76 & 22.72 & 21.32 \\
        No Transformers & 27.74 & 27.21 & 26.64 & 26.03 & 25.42 & 24.74 & 24.01 & 23.19 & 22.21 & 20.88 \\
        No U-Net & 28.33 & 27.73 & 27.14 & 26.49 & 25.79 & 25.04 & 24.26 & 23.35 & 22.33 & 20.99 \\
        \underline{Baseline} & \underline{29.87} & \underline{29.27} & \underline{28.64} & \underline{27.97} & \underline{27.25} & \underline{26.49} & \underline{25.62} & \underline{24.65} & \underline{23.47} & \underline{21.90} \\
        \bottomrule
    \end{tabular}
    \caption{Ablation Study Results - PSNR}
    \label{table:ablation-2}
\end{table}

\begin{table}[!hbt]
    \centering
    \begin{tabular}{>{\centering\arraybackslash}p{2.5cm}*{10}{>{\centering\arraybackslash}p{.55cm}}}
        \toprule
        Noise Density & 150 & 160 & 170 & 180 & 190 & 200 & 210 & 220 & 230 & 240 \\
        \midrule
        \textbf{BERNOULLI} \\
        5 Transform & 0.926 & 0.916 & 0.904 & 0.890 & 0.873 & 0.853 & 0.826 & 0.791 & 0.740 & 0.661 \\
        No Anisotropic Block & 0.925 & 0.916 & 0.904 & 0.889 & 0.873 & 0.852 & 0.826 & 0.790 & 0.740 & 0.661 \\
        No SeConv & 0.921 & 0.909 & 0.897 & 0.882 & 0.864 & 0.842 & 0.814 & 0.778 & 0.725 & 0.645 \\
        No Transform & 0.906 & 0.894 & 0.880 & 0.864 & 0.844 & 0.820 & 0.790 & 0.751 & 0.699 & 0.619 \\
        No U-Net & 0.911 & 0.899 & 0.886 & 0.869 & 0.850 & 0.826 & 0.796 & 0.757 & 0.703 & 0.623 \\
        \underline{Baseline} & \underline{0.930} & \underline{0.921} & \underline{0.909} & \underline{0.896} & \underline{0.880} & \underline{0.860} & \underline{0.835} & \underline{0.800} & \underline{0.750} & \underline{0.671} \\
        \midrule
        \textbf{POISSON} \\
        5 Transform & 0.939 & 0.936 & 0.933 & 0.929 & 0.926 & 0.921 & 0.918 & 0.916 & 0.912 & 0.908 \\
        No Anisotropic Block & 0.941 & 0.938 & 0.935 & 0.930 & 0.927 & 0.924 & 0.919 & 0.916 & 0.912 & 0.908 \\
        No SeConv & 0.941 & 0.937 & 0.935 & 0.931 & 0.927 & 0.924 & 0.919 & 0.916 & 0.912 & 0.907 \\
        No Transform & 0.920 & 0.915 & 0.911 & 0.906 & 0.901 & 0.896 & 0.890 & 0.884 & 0.878 & 0.869 \\
        No U-Net & 0.935 & 0.930 & 0.927 & 0.923 & 0.919 & 0.915 & 0.911 & 0.906 & 0.902 & 0.897 \\
        \underline{Baseline} & \underline{0.949} & \underline{0.946} & \underline{0.944} & \underline{0.941} & \underline{0.937} & \underline{0.933} & \underline{0.930} & \underline{0.926} & \underline{0.923} & \underline{0.917} \\
        \midrule
        \textbf{SAP} \\
        5 Transform & 0.922 & 0.912 & 0.900 & 0.886 & 0.869 & 0.848 & 0.820 & 0.784 & 0.733 & 0.652 \\
        No Anisotropic Block & 0.923 & 0.913 & 0.901 & 0.886 & 0.869 & 0.848 & 0.821 & 0.784 & 0.733 & 0.652 \\
        No SeConv & 0.917 & 0.906 & 0.893 & 0.878 & 0.859 & 0.836 & 0.808 & 0.769 & 0.717 & 0.636 \\
        No Transform & 0.897 & 0.885 & 0.870 & 0.853 & 0.833 & 0.807 & 0.776 & 0.737 & 0.682 & 0.601 \\
        No U-Net & 0.908 & 0.896 & 0.883 & 0.867 & 0.847 & 0.823 & 0.792 & 0.752 & 0.696 & 0.609 \\
        \underline{Baseline} & \underline{0.929} & \underline{0.919} & \underline{0.909} & \underline{0.895} & \underline{0.879} & \underline{0.859} & \underline{0.832} & \underline{0.797} & \underline{0.747} & \underline{0.666} \\
        \bottomrule
    \end{tabular}
    \caption{Ablation Study Results - SSIM}
    \label{table:ablation-3}
\end{table}

The results presented in Tables \ref{table:ablation-1}, \ref{table:ablation-2}, and \ref{table:ablation-3} validate the effectiveness of MPDenoiseNet's multi-component design. Removing any major architectural element consistently leads to a degradation in performance across all tested noise types and metrics compared to the full baseline model, indicating that each module positively contributes to the final denoising capability. The magnitude of the performance drop upon removal allows us to rank the components by their impact. The Transformer blocks ($f_{TF}$) exhibit the most significant influence, as their removal causes the largest decrease in both PSNR and SSIM values across all noise conditions, highlighting the importance of global context modeling provided by the attention mechanism. The asymmetric U-Net AutoEncoder ($f_{AE}$) demonstrates the second highest impact, suggesting its crucial role in refining features and reconstructing the final high-quality image. The SeConv blocks ($f_{SeConv}$) show a notable contribution, particularly evident in the performance drop for Bernoulli and Salt-and-Pepper noise (Table \ref{table:ablation-2}), validating their specific design goal for impulse noise removal. Finally, the Anisotropic Diffusion blocks ($f_{AD}$), while contributing positively, show the smallest relative impact among the removed components in this study. These findings confirm the synergistic contribution of MPDenoiseNets architectural elements and justify the inclusion of each specialized component.